\documentclass[12pt, a4paper]{article}

\usepackage[utf8]{inputenc}
\usepackage[T1]{fontenc}
\usepackage[russian]{babel}
\usepackage[oglav,spisok,boldsect,eqwhole,figwhole,hyperref,hyperprint,remarks,greekit]{./style/fn2kursstyle}
\graphicspath{{./style/}{./figures/}}

\usepackage{multirow}
\usepackage{supertabular}
\usepackage{multicol}
\usepackage{hyperref}
\usepackage{enumitem}
% Параметры титульного листа
\title{Лабораторная работа 1}
\author{А.\,Р.~Ластра-Грек}
\supervisor{А.\,В.~Череднеченко}
\group{ФН2-51Б}
\date{2023}

% Переопределение команды \vec, чтобы векторы печатались полужирным курсивом
\renewcommand{\vec}[1]{\text{\mathversion{bold}${#1}$}}%{\bi{#1}}
\newcommand\thh[1]{\text{\mathversion{bold}${#1}$}}
%Переопределение команды нумерации перечней: точки заменяются на скобки
\renewcommand{\labelenumi}{\theenumi)}

\usepackage{amsmath}
\usepackage{amssymb}

\usepackage{float, graphicx}

\usepackage{nicefrac}

\usepackage{relsize}
\usepackage{caption}
\usepackage{listings}
\usepackage{xcolor}
\lstset { %
	language=C++,
	%backgroundcolor=\color{black!5}, % set backgroundcolor
	basicstyle=\footnotesize % basic font setting
}

\frenchspacing
\righthyphenmin=2
\linespread{1.25}

\newcounter{controlQCounter}
\setcounter{controlQCounter}{1}

\DeclareMathOperator{\cond}{cond}
\newcommand{\tr}{\ensuremath{\mathrm{T}}}
\newcommand{\bigO}{\ensuremath{\mathrm{O}}}
\newcommand{\linSpan}[2]{\ensuremath{\overline{#1, \, #2}}}
\newcommand{\RR}{\mathbb{R}}
\newcommand{\NN}{\mathbb{N}}

\newcommand{\clock}[4]{\boldsymbol{\mathrm{#1}\!:\!\mathrm{#2} \to \mathrm{#3}\!:\!\mathrm{#4} \blacktriangleright} \quad}
\newcommand{\point}{$\_ \_ \_ \_ \_ \_ \_ \_ \_ \_ \_ \_ \_ \_ \_ \_ \_ \_ \_ \_ \_ \_ \_ \_ \_ $}
\newcommand{\size}{1}

\newenvironment{controlQ}{\textbf{Вопрос \arabic{controlQCounter}.} \addtocounter{controlQCounter}{1} \itshape}{\upshape}
\newenvironment{controlAns}{\textrm{Ответ:} \par }{ \bigskip }

\begin{document}
	%\maketitle
	
	\tableofcontents
	
	\newpage
	
	\section{Введение}
	\subsection{Мотивация}
	
	\newpage
	
	\section{Семейство методов Курганова-Тадмора}
	
	В данной работе будет рассмотрено три схемы, реализующие семейство методов Курганова-Тадмора: схема $SD2$ (Semi-Discrete $2$nd order), схема $SD3$ (Semi-Discrete $3$nd order) и $FD2$ (Fully-Discrete $2$nd order).
	
	\subsection{Постановка задачи и метод конечных объемов}
	
	Рассмотрим одномерное уравнение:
	
	\[
	\frac{\partial u}{\partial t} + \frac{\partial f(u)}{\partial x} = 0.
	\]
	
	Методы Курганова-Тадмора являются консервативными, для них характерно понятие ячейки.
	Разобьем пространство на ячейки шириной $\Delta x$. Обозначим $x_i = i \Delta x$.
	В методе конечных объемов мы ищем эволюцию средних значений по ячейке $\bar{u}_i(t)$
	\[
	\bar{u}_i(t) = \frac{1}{\Delta x} \int_{x_{i-1/2}}^{x_{i+1/2}} u(\xi, t) \, d\xi
	\]
	
	(Так как я еще нигде не описывал метод Рунге-Кутты, наверное здесь следует написать о том что эти методы находят лишь правую часть уравнения, то есть они работают по пространству)
	
	
	\subsection{Метод SD2 и SD3}
	
	Рассмотрим основную идею метода SD2. При разбиении пространства на ячейки мы знаем лишь среднее значение плотности воды или газа внутри. Если происходит ударная волна при поиске решения в виде сеточной функции (дописать определение) может возникнуть размывание и дрожание решения. Однако если мы будем учитывать не только среднее значении внутри, но и значения с каждой стороны от стыка ячейки, можно получить более точное решение. Первым шагом в методе SD2 происходит реконструкция, то есть поиск значений на краях ячейки. Для этого мы находим разностные производные вперед и назад
	
	\[
		f'(x_i) \approx \frac{f(x_{i+1}) - f(x_i)}{h} ,\quad	f'(x_i) \approx \frac{f(x_i) - f(x_{i-1})}{h} 
	\]
	
	после, используя лимитер следующего вида находим значения наклонов в каждой ячейке: (наверное стоит добавить определение того что такое лимитер)
	
	\[
		\text{minmod}(a,b) = \frac{1}{2} \big( \operatorname{sgn}(a) + \operatorname{sgn}(b) \big) \cdot \min\big( |a|, |b| \big),
	\]
	
	где \( \operatorname{sgn}(x) \) — знаковая функция:
	\[
	\operatorname{sgn}(x) =
	\begin{cases} 
		1, & x > 0, \\
		0, & x = 0, \\
		-1, & x < 0,
	\end{cases}
	\]
	
	\( |x| \) — модуль числа \( x \), \( \min(|a|, |b|) \) — минимум из \( |a| \) и \( |b| \).
	
	$u_i^{Left} = u_i - S_i \frac{\Delta x}{2}$ и $u_i^{Right} = u_i + S_i \frac{\Delta x}{2}$
	
	\subsection{Метод FD2}
	
	\section{Программная реализация и анализ результатов}
	\subsection{Разработка и реализация вычислительной библиотеки на основе JAX}
	\subsection{Верификация численного метода и оценка точности на тестовых задачах}
	\subsection{Анализ производительности и сравнение версий}
	
	\section{Заключение}
	\section{Список литературы}
	
\end{document}